(1) 在近轴条件下，由不同入射方向的波前倾斜产生两缝入射相位差 $-kd\alpha$，而向观察方向 $\theta$ 传播产生的相位差为 $kd\theta$，其中 $k=2\pi/\lambda$。因此对给定入射角，两缝同强度的远场干涉强度可写为
\[
I_\alpha(\theta)=2I_0p(\alpha)[1+\cos(kd(\theta-\alpha))],
\]
这里 $I_0$ 定义为全部角分布积分后任一单缝的单独光强，且两缝理想窄、忽略单缝衍射包络变化。

(2) 由于不同入射角属于相互不相干的光源分量，实验探测的总强度是 $\int I_\alpha\,d\alpha$，而不是各角复振幅相加。利用高斯分布的特征函数
\[
\int_{-\infty}^\infty p(\alpha)e^{-ikd\alpha}\,d\alpha
=\exp[-(kd\sigma)^2/2]
\]
且角分布关于零对称，可得
\[
I(\theta)=2I_0\left[1+
e^{-(kd\sigma)^2/2}\cos(kd\theta)\right].
\]
条纹的最大、最小强度为 $2I_0(1\pm V)$，故可见度
\[
V=\frac{I_{\max}-I_{\min}}{I_{\max}+I_{\min}}
=e^{-(kd\sigma)^2/2}
=\exp[-2\pi^2d^2\sigma^2/\lambda^2].
\]

(3) 定义 $V(d_{\rm coh})=1/e$，则 $(kd_{\rm coh}\sigma)^2/2=1$，所以
\[
d_{\rm coh}=\frac{\sqrt2}{k\sigma}
=\frac{\lambda}{\sqrt2\pi\sigma}.
\]
光源角宽越大，横向相干长度越短；波长越长，在给定角尺度下的相干长度越长。此长度定义依赖于指定的高斯模型和 $1/e$ 阈值。

(4) 当 $\sigma\to0^+$ 时光源趋于单一方向，$V\to1$，恢复完全相干双缝干涉。固定 $\sigma>0$ 而 $d\to\infty$，$V\to0$，条纹被各入射方向相位平均抹掉，但单缝总强度仍存在。若错误地把不同 $\alpha$ 的复振幅先相加再取模平方，就会引入源分量之间实际上不存在的干涉项，得到错误的相干程度。这一处理体现了光源角谱与横向相干度的傅里叶关系。
